\documentclass[11pt,a4paper]{amsart}

\pdfinfoomitdate=1
\pdftrailerid{}
\pdfsuppressptexinfo=15

\usepackage[T1]{fontenc}
\usepackage{lmodern}
\usepackage[margin=27mm]{geometry}
\usepackage{microtype}
\usepackage{amsmath,amssymb,mathtools}
\usepackage{booktabs}
\usepackage{xcolor}
\usepackage[numbers,sort&compress]{natbib}
\usepackage[colorlinks=true,linkcolor=blue!55!black,citecolor=green!45!black,urlcolor=blue!60!black]{hyperref}
\usepackage[nameinlink,noabbrev]{cleveref}

\raggedbottom

\newcommand{\Z}{\mathbb Z}
\newcommand{\I}{\mathcal I}
\newcommand{\Ann}{\operatorname{Ann}}
\newcommand{\Fix}{\operatorname{Fix}}
\newcommand{\Rec}{\operatorname{Rec}}
\newcommand{\depth}{\operatorname{depth}}

\newtheorem{theorem}{Theorem}[section]
\newtheorem{proposition}[theorem]{Proposition}
\newtheorem{corollary}[theorem]{Corollary}
\newtheorem{remark}[theorem]{Remark}

\title[Annihilator--power ideal dynamics]{Annihilator--Power Dynamics on Ideals of Residue Rings:\ Resonant Fixed Points, Two-Cycles, and Exact Transients}
\author{Anonymous}
\date{Internal Stage 2 draft, 29 August 2026}

\begin{document}

\begin{abstract}
Fix $r\geq2$.  On the ideal lattice of $\Z/N\Z$ we study
$T_r(I)=\Ann(I)^r$.  For a prime-power component $p^a$, the valuation
coordinate is the clipped reflection
\[
 f_{a,r}(e)=\min\{a,r(a-e)\},\qquad 0\leq e\leq a.
\]
The endpoints form a two-cycle, while an interior fixed point exists exactly
when $r+1$ divides $a$.  Before clipping, the centered deviation
$\Delta(e)=(r+1)e-ra$ is multiplied by $-r$ at every step.  This gives an
exact hitting-time formula, a closed transient cumulative distribution, and
a sharp logarithmic maximum depth.  Chinese remaindering then yields the
complete recurrent set, all fixed and two-cycle counts, and the Artin--Mazur
zeta function for arbitrary $N$.  Literal ideal arithmetic for every
$N\leq1000$ independently checks the valuation proof.
\end{abstract}

\maketitle

\section{Map and ownership boundary}

Let $N\geq2$, $R_N=\Z/N\Z$, and let $\I_N$ be its finite set of ideals.
For a fixed integer $r\geq2$, define
\begin{equation}
 \label{eq:map}
 T_r\colon\I_N\longrightarrow\I_N,
 \qquad T_r(I)=\Ann_{R_N}(I)^r.
\end{equation}
Ideal annihilators, ideal powers, and the Chinese remainder theorem are
standard commutative-algebra ingredients; see \citet{AtiyahMacdonald1969}.
Annihilating-ideal graphs and their residue-ring specializations are also an
established topic \citep{BehboodiRakeei2011,AfkhamiEtAl2016}.  We assign
these ingredients no novelty credit.  More directly, Schwiebert treats the
radical and annihilator as operators on the ideal poset, studies their
generated monoid, and analyzes finite-product behavior
\citep{Schwiebert2017}.  That operator-dynamics viewpoint, the bare
annihilator iteration, and its product-decomposition framework are owned
background here.  Our bounded target is only the temporal package for the
specific power-after-annihilator self-map \eqref{eq:map}: the
clipped-reflection normal form, the divisibility-controlled fixed-point
anomaly, exact transient layers, and the full one/two-cycle census.  A
targeted search completed on 29 August 2026 did not locate that conjunction.
This is not an exhaustive owner certificate; external circulation and every
priority claim remain on \textbf{HOLD}.

Factor $N$ as
\begin{equation}
 \label{eq:factor}
 N=\prod_{i=1}^{s}p_i^{a_i}.
\end{equation}
The Chinese remainder theorem identifies $\I_N$ with
$\prod_i\{0,1,\ldots,a_i\}$: the coordinate $e_i$ denotes the ideal
$(p_i^{e_i})$ in $\Z/p_i^{a_i}\Z$.  The first proof route works entirely in
these valuation coordinates.  The control route instead represents an ideal
of $R_N$ by its unique divisor generator $(d)$, computes
\begin{equation}
 \label{eq:literal}
 T_r((d))=\left(\gcd\!\left(N,(N/d)^r\right)\right),
\end{equation}
and only then compares with the coordinate formulas.

\section{Prime-power normal form}

Fix $a\geq1$ and abbreviate $f=f_{a,r}$.  Since
$\Ann((p^e))=(p^{a-e})$ and $(p^b)^r=(p^{\min(a,rb)})$, we have the following
exact model.

\begin{proposition}[Clipped reflection]
\label{prop:normal}
On $\I_{p^a}$, the map \eqref{eq:map} is conjugate to
\begin{equation}
 \label{eq:f}
 f(e)=\min\{a,r(a-e)\},\qquad 0\leq e\leq a.
\end{equation}
Put $\Delta(e)=(r+1)e-ra$.  Whenever $f(e)<a$,
\begin{equation}
 \label{eq:deviation}
 \Delta(f(e))=-r\Delta(e).
\end{equation}
Clipping occurs precisely when $\Delta(e)\leq-a/r$.
\end{proposition}

\begin{proof}
Only the deviation statements require verification.  In the unclipped
branch $f(e)=r(a-e)$, whence
\[
 (r+1)f(e)-ra
 =r\bigl((r+1)(a-e)-a\bigr)=-r\Delta(e).
\]
The inequality $r(a-e)\geq a$ is equivalent to
$(r+1)e-ra\leq-a/r$.
\end{proof}

The endpoints satisfy $0\leftrightarrow a$.  An interior point is fixed if
and only if $e=r(a-e)$, which already exposes the arithmetic anomaly.

\begin{theorem}[Recurrent states and resonance]
\label{thm:recurrent-coordinate}
The recurrent set of $f_{a,r}$ is
\begin{equation}
 \label{eq:rec-coordinate}
 \{0,a\}\ \cup\
 \left\{\frac{ra}{r+1}:\ (r+1)\mid a\right\}.
\end{equation}
The endpoints form one two-cycle.  The displayed interior state, when it
exists, is fixed.  There are no other cycles.
\end{theorem}

\begin{proof}
The endpoint statements and the fixed-point equation follow from
\eqref{eq:f}.  Suppose $\Delta(e)\neq0$ and the orbit has not clipped.  By
\eqref{eq:deviation}, its deviation alternates sign and its absolute value is
multiplied by $r>1$.  It therefore eventually enters the clipping region.
The next state is $a$, hence belongs to the endpoint two-cycle.  Thus every
nonfixed interior state is transient.
\end{proof}

\section{Exact transient layers}

For a state $e$, let $\depth(e)$ be the first time its orbit is recurrent.
The deviation gives a closed hitting-time clock rather than only an upper
bound.

\begin{theorem}[Exact depth]
\label{thm:depth-coordinate}
For $e$ in \eqref{eq:rec-coordinate}, $\depth(e)=0$.  For every other state,
put $D=\Delta(e)$.  If $D<0$, let
\[
 k_- =\min\{k\geq0:r^{2k+1}|D|\geq a\};
\]
if $D>0$, let
\[
 k_+ =\min\{k\geq0:r^{2k+2}D\geq a\}.
\]
Then
\begin{equation}
 \label{eq:depth-coordinate}
 \depth(e)=
 \begin{cases}
  2k_-+1,&D<0,\\
  2k_++2,&D>0.
 \end{cases}
\end{equation}
Consequently, on each sign side that contains transient states, the largest
depth occurs at the smallest available value of $|D|$.  The global maximum
is the larger of those sidewise values, or zero when there are no transient
states.  In particular it is $O(\log_r a)$, and
\eqref{eq:depth-coordinate} is a sharp formula for every $(a,r)$, including
$a=1$.
\end{theorem}

\begin{proof}
Before clipping, the deviation at time $j$ is $(-r)^jD$.  A negative
deviation clips on the following update exactly when its magnitude times
$r$ is at least $a$.  If $D<0$, negative signs occur at times $2k$; if
$D>0$, they occur at times $2k+1$.  The two displayed threshold conditions
and the additional clipping update give \eqref{eq:depth-coordinate}.
Monotonicity in $|D|$ proves the sharp maximum statement.
\end{proof}

For completeness, the whole transient distribution also has a floor/ceiling
form.  Put $\varepsilon_{a,r}=\mathbf1_{(r+1)\mid a}$ and, for $L\geq1$,
\begin{align}
 M_-(L)&=\max\left\{0,
  \min\left(a-1,\left\lfloor\frac{ra-L}{r+1}\right\rfloor\right)\right\},
 \label{eq:Mminus}\\
 M_+(L)&=\max\left\{0,
  a-\max\left(1,\left\lceil\frac{ra+L}{r+1}\right\rceil\right)\right\}.
 \label{eq:Mplus}
\end{align}
These count interior states with $\Delta\leq-L$ and $\Delta\geq L$,
respectively.

\begin{corollary}[Transient cumulative distribution]
\label{cor:cdf}
Let $C_{a,r}(t)=\#\{e:\depth(e)\leq t\}$.  Then
$C_{a,r}(0)=2+\varepsilon_{a,r}$.  For $t\geq1$, set
\[
 j_-(t)=2\left\lfloor\frac{t-1}{2}\right\rfloor,
 \qquad L_-(t)=\left\lceil\frac{a}{r^{j_-(t)+1}}\right\rceil.
\]
For $t\geq2$, also set
\[
 j_+(t)=2\left\lfloor\frac{t-2}{2}\right\rfloor+1,
 \qquad L_+(t)=\left\lceil\frac{a}{r^{j_+(t)+1}}\right\rceil.
\]
Then
\begin{equation}
 \label{eq:cdf}
 C_{a,r}(t)=2+\varepsilon_{a,r}+M_-(L_-(t))
 +\mathbf1_{t\geq2}M_+(L_+(t)).
\end{equation}
Thus the exact depth-$t$ shell is $C_{a,r}(t)-C_{a,r}(t-1)$.
\end{corollary}

\begin{proof}
Among negative starting deviations, the largest even exponent available by
time $t-1$ is $j_-(t)$; among positive ones, the largest admissible odd
exponent is $j_+(t)$, which exists only for $t\geq2$.  The clipping test in
\cref{prop:normal} becomes $|D|\geq L_-(t)$ or $D\geq L_+(t)$.  Counting the
arithmetic progression $D=(r+1)e-ra$ over $1\leq e\leq a-1$ gives
\eqref{eq:Mminus}--\eqref{eq:Mplus}.  The endpoints and possible resonant
fixed point contribute $2+\varepsilon_{a,r}$.
\end{proof}

\section{Chinese products, cycles, and zeta}

Return to \eqref{eq:factor} and put
\begin{equation}
 \label{eq:AB}
 A=\prod_{i=1}^{s}\varepsilon_{a_i,r},
 \qquad B=\prod_{i=1}^{s}(2+\varepsilon_{a_i,r}).
\end{equation}

\begin{theorem}[Complete finite dynamics]
\label{thm:global}
Under Chinese remaindering, $T_r$ is the product of the maps
$f_{a_i,r}$.  Its recurrent set has size $B$, its maximum transient depth is
\begin{equation}
 \label{eq:global-depth}
 \max_{1\leq i\leq s}\ \max_{0\leq e\leq a_i}\depth_{a_i,r}(e),
\end{equation}
and the number of states of depth at most $t$ is
\begin{equation}
 \label{eq:global-cdf}
 \prod_{i=1}^{s}C_{a_i,r}(t).
\end{equation}
For every $k\geq1$,
\begin{equation}
 \label{eq:fixed-global}
 \#\Fix(T_r^k)=
 \begin{cases}
  A,&k\text{ odd},\\
  B,&k\text{ even}.
 \end{cases}
\end{equation}
Hence there are $A$ fixed points and $(B-A)/2$ two-cycles, and the
Artin--Mazur zeta function \citep{ArtinMazur1965} is
\begin{equation}
 \label{eq:zeta}
 \zeta_{T_r}(z)=(1-z)^{-A}(1-z^2)^{-(B-A)/2}.
\end{equation}
No period greater than two occurs.
\end{theorem}

\begin{proof}
Annihilators, products, and powers commute with the finite product
decomposition of $R_N$, giving the coordinate product.  A product point is
recurrent exactly when each coordinate is recurrent; every such coordinate
has period one or two, so the product period also divides two.  Depth is the
maximum coordinate depth, proving \eqref{eq:global-depth}, while independent
coordinate choices give \eqref{eq:global-cdf}.  An odd iterate fixes only
the resonant fixed coordinate in every component, producing $A$.  An even
iterate fixes every recurrent coordinate, producing $B$.  Subtracting the
fixed points from the points fixed by the square and dividing by two gives
the cycle counts and \eqref{eq:zeta}.
\end{proof}

The dynamics depends on the prime exponents and not on the primes themselves:
moduli with the same exponent multiset have conjugate systems.  This is a
nonrigidity statement, not a claim that the phase portrait always recovers
that multiset.

The internal collision boundary is also explicit.  P100 uses least-valuation
digit erasure on residue classes and reconstructs prime data from a transient
profile; it does not act on the ideal lattice or alternate an annihilator with
an ideal power.  P102 uses an involution--norm map in a cyclic group algebra;
its orbit structure comes from polynomial factorization rather than the
clipped expanding reflection \eqref{eq:f}.  Shared valuation, CRT, and zeta
bookkeeping receive no contribution credit.

\section{Independent exact control}

The accompanying standard-library verifier implements two independent
lanes.  The first enumerates $e=0,\ldots,a$, checks the deviation clock,
the cumulative formula, and all iterate-fixed counts.  The second enumerates
actual divisor ideals $(d)$ for every $N\leq1000$, applies the literal gcd
rule \eqref{eq:literal}, and compares trajectories, depths, recurrent sets,
and fixed counts with the CRT prediction.  These finite computations are
falsification controls, not proofs of the quantified theorems.

\section{Conclusion}

Annihilation followed by an ideal power becomes a clipped expanding
reflection on each prime exponent.  Its sign alternation forces a universal
two-cycle, while the single equation $(r+1)e=ra$ creates a sharp divisibility
resonance.  The same normal form supplies every transient layer and, after
Chinese remaindering, the complete cycle and zeta data.  Specialist owner
review remains necessary before any external use.

\bibliographystyle{plainnat}
\bibliography{references}

\end{document}
